\begin{center}\begin{tikzpicture}[font=\small,>=stealth,line join=round]
\draw[gray!40,rounded corners=3pt] (0,0) rectangle (14.2,8.3);\node at (7.1,7.700000000000001) {每次减去的数若改变，原值就未必还原};
\begin{scope}[shift={(2.6,2.8)},x=1.0cm,y=0.7666666666666666cm]\draw[->,gray] (-1.5,0)--(4.5,0) node[right]{$x$};\draw[->,gray] (0,-1.5)--(0,4.5) node[above]{$y$};\draw[blue,thick,solid,domain=-1:4,samples=180] plot(\x,{\x});\fill[orange!85!black] (0,0) circle(2.2pt);\fill[orange!85!black] (1,1) circle(2.2pt);\fill[orange!85!black] (2,2) circle(2.2pt);\end{scope}\node[anchor=west,align=left,text width=5.0cm] at (8.4,4.0) {$f(x)=x$\\[9pt]令 $c(x)=2x+1$\\[9pt]则 $f(x+1)=c(x)-f(x)$\\[9pt]却有 $f(x+2)\ne f(x)$\\[9pt]这不是固定常数 $c$};
\node[text width=12.6cm,align=center] at (7.1,.65) {两次操作必须使用同一个规则与同一个常数；否则抵消过程不成立。};
\end{tikzpicture}\end{center}
